\chapter*{Resumen}
La variabilidad entre personas limita la generalización de las interfaces cerebro-computador basadas en imaginación motora y dificulta comparar arquitecturas evaluadas con protocolos diferentes. Se evaluó la clasificación de imaginación de movimiento de las manos izquierda y derecha con electroencefalografía superficial en OpenBMI, conjunto público de 54 participantes, y se empleó WBCIC-SHU para validación externa con 45 canales comunes. Las señales se filtraron, limpiaron, segmentaron y normalizaron antes de su adaptación a cuatro arquitecturas. EEGNet sirvió como referencia; Conformer y EEGMamba se entrenaron desde cero, mientras LaBraM Base se ajustó completamente desde pesos preentrenados. Se aplicaron evaluación intra-sujeto por sesión y validación dejando un sujeto fuera (\loso{}), con las semillas 2026, 2027 y 2028. Un indicador comparó personas vistas y no vistas con los mismos pesos para separar ese cambio del volumen de entrenamiento; la validación externa examinó además el cambio de conjunto y protocolo. En \loso{}, las exactitudes medias sobre los 54 sujetos y tres semillas fueron 0.761 para EEGNet, 0.747 para Conformer, 0.755 para LaBraM y 0.715 para EEGMamba. Los regímenes de optimización no fueron idénticos entre todas las semillas. EEGNet presentó el menor indicador de degradación entre personas y EEGMamba el menor indicador de cambio de protocolo. La evaluación incluyó incertidumbre, estabilidad y criterios adicionales para una decisión comparativa. La regla, fijada después del primer análisis para el conjunto principal y los regímenes evaluados, prefirió EEGNet por su menor costo medido entre las tres arquitecturas equivalentes en exactitud \loso{} y degradación según la prueba bayesiana y TOST con Holm. Esta preferencia no implica superioridad en todos los ejes ni fuera de las condiciones evaluadas. Los resultados delimitan la interpretación de la generalización a los conjuntos y condiciones estudiados y ofrecen una base reproducible para comparar modelos y orientar investigaciones posteriores.

\noindent\textbf{Palabras clave:} interfaz cerebro-computador; imaginación motora; electroencefalografía; cambio de dominio; aprendizaje profundo; generalización.
